%%=============================================================================
%%    Conclusie
%%=============================================================================

%    Hierin ga je na welke technische doelstellingen bereikt zijn. Vertrekkende van de bachelorproefopdracht ga
%    je na welke technische doelstellingen al dan niet bereikt zijn. Probeer de waarde van je bachelorproef in te
%    schatten. Volgende vragen kunnen hierbij helpen:
%    ● Welke doelstellingen zijn bereikt?
%    ● Als er bepaalde doelstellingen niet bereikt zijn, wat zijn de redenen hiervoor? Hoe kon dit
%       voorkomen worden?
%    ● Gaan jouw bereikte resultaten in de toekomst toegepast worden? Heb je daar stappen voor
%       ondernomen?
%    ● Hoe kunnen je resultaten in de toekomst nog verbeterd worden?
%    ● Is er nog vervolgonderzoek nodig?
%    ● Misschien heb je ervaren dat bepaalde werksituaties minder efficiënt aangepakt worden. Heb jij tips
%       om deze op te lossen?

\chapter*{Conclusie}%
\label{ch:conclusie}


Hierin ga je na welke technische doelstellingen bereikt zijn. Vertrekkende van de bachelorproefopdracht ga
je na welke technische doelstellingen al dan niet bereikt zijn. Probeer de waarde van je bachelorproef in te
schatten. Volgende vragen kunnen hierbij helpen:

\begin{itemize}[]
    \item Welke doelstellingen zijn bereikt?
    \item Als er bepaalde doelstellingen niet bereikt zijn, wat zijn de redenen hiervoor? Hoe kon dit voorkomen worden?
    \item Gaan jouw bereikte resultaten in de toekomst toegepast worden? Heb je daar stappen voor ondernomen?
    \item Hoe kunnen je resultaten in de toekomst nog verbeterd worden?
    \item Is er nog vervolgonderzoek nodig?
    \item Misschien heb je ervaren dat bepaalde werksituaties minder efficiënt aangepakt worden. Heb jij tips om deze op te lossen?
\end{itemize}
